\section{ĐỀ SỐ \thesection } %star8
\setcounter{bt}{0}

%%%%%%%%%%%%%%% CÂU 1
\begin{bt} Tính hợp lý
	\begin{listEX}[2]
		\item $2\cdot 3^{2}-32:2^{2}$;
		\item $19 - 950 : \{ 2 \cdot [ 51 - (15 - 2^{2} \cdot 3)^{3} ] \}$;
		\item $(-3)\cdot 5\cdot (-2):(-15)$;
		\item $63\cdot 57+63\cdot 43-300$.
	\end{listEX}
	\loigiai{ a) $6$ \hfill b) $0$ \hfill c) $-2$ \hfill d) $6\,000$
	}
\end{bt}

%%%%%%%%%%%%%%% BÀI 2
\begin{bt}  Tìm $x \in \mathbb{Z}$, biết
	\begin{listEX}[2]
		\item $216-(72+x)=48$;
		\item $(x-24)^{5}\cdot 7^{8}=7^{18}$;
		\item $3\cdot\left|-x-5\right|=(-3)^{2}$;
		\item $2^{2x+4}\cdot 2^{3x}-2^{x}\cdot 2^{4x+2}=384$.
	\end{listEX}
	\loigiai{a) $x=96$ \hfill b) $x=73$ \hfill c) $x=2$ hay $x=8$ \hfill d) $x=1$}
\end{bt}

%%%%%%%%%%%%%%% BÀI 3
\begin{bt}  Số học sinh khối $6$ của một trường THCS khi xếp hàng $12$, $15$, $18$ thì đều thừa $4$ học sinh, nhưng nếu xếp hàng $26$ thì vừa đủ. Tính số học sinh khối $6$ của trường THCS đó, biết số học sinh khối $6$ nhỏ hơn $400$ em.
	\loigiai{$364$}
\end{bt}

%%%%%%%%%%%%%%% BÀI 4
\begin{bt} 
	Một nhà sách nhập về một lô hàng gồm $40$ cái máy tính Casio–FX580 với giá $14\,000\,000$ đồng và bán ra với giá $400\,000$ đồng một máy. Sau khi bán được $20$ cái máy tính thì nhà sách giảm giá còn $380\,000$ đồng/máy và bán hết số máy tính còn lại. Hỏi nhà sách lời hay lỗ bao nhiêu tiền khi bán hết lô hàng trên?
	\loigiai{
		Nhà sách lời $1\,600\,000$
	}
\end{bt}

%%%%%%%%%%%%%%% BÀI 5
\begin{bt}%[Bài toán thực tế hình học]
	Một mảnh đất hình bình hành $ABCD$ có cạnh đáy $AD = 6$ m và chiều cao tương ứng $AE = 3$ m.  
	\begin{center}
		\begin{tikzpicture}[scale=1,line width=0.8pt,line join=round,>=stealth]
		 
			\def\b{3}  
			\def\t{1.2}  
			 
			\path 
			(0,0) coordinate (B)
			(0:2*\b) coordinate (C)
			(0:0.7*\b)coordinate(E)
			(E)++(90:\b)coordinate(A)
			(A)++(0:2*\b)coordinate(D)
			(E)++(0:\t)coordinate(F)
			(F)++(90:\b)coordinate(G)
			;
			
		 
			\draw (A)--(B)--(C)--(D)--cycle;
		 
			\fill[pattern=north east lines, pattern color=blue](A)--(E)--(F)--(G);
			 \draw   (A)--(E)--(F)--(G);
		  %%%% Hiển thị tên điểm
		 \foreach \x/\g in {A/90,B/-90, C/-90,D/90,E/-90,F/-90,G/90
		 }\draw[fill=white](\x)circle(0pt)+(\g:.3)node{$\x$};
		 \draw pic[draw, angle radius=2.5mm]{right angle=A--E--B};
		\end{tikzpicture}
	\end{center}
	
	\begin{listEX}
		\item Tính diện tích mảnh đất hình bình hành đó.
		\item Chủ mảnh đất làm một lối đi hình chữ nhật $AEFG$ có kích thước $AG = 12$ dm, $AE = 3$ m.  
		Hỏi diện tích mảnh đất hình bình hành gấp mấy lần diện tích lối đi hình chữ nhật?
		\item Để lối đi sạch đẹp hơn, chủ nhà lát những viên gạch hình vuông kín lối đi.  
		Hỏi chủ nhà sử dụng những viên gạch hình vuông có độ dài cạnh lớn nhất bằng bao nhiêu để không có viên gạch nào bị cắt xén (coi mạch vữa không đáng kể) và khi đó người chủ cần bao nhiêu viên gạch?
	\end{listEX}
	
	\loigiai{a) $18$ m$^2$ \hfill b) $5$ (lần) \hfill c) $10$ (viên)
	}
\end{bt}


%%%%%%%%%%%%%%% BÀI 6
\begin{bt}%[Thống kê huy chương theo điểm]
	Trong một kỳ thi quốc tế, các thành viên của đội tuyển Việt Nam đạt được số điểm lần lượt là:
	$26;  18; 36;  28;  21; 32.$
	Biết rằng người ta phân loại huy chương theo điểm số như sau:
	 \begin{center}
	\begin{tabular}{|c|c|}
		\hline
		\textbf{Khoảng điểm} & \textbf{Loại huy chương} \\ \hline
		Từ $15$ đến dưới $24$ & Huy chương đồng \\ \hline
		Từ $24$ đến dưới $31$ & Huy chương bạc \\ \hline
		Từ $31$ trở lên & Huy chương vàng \\ \hline
	\end{tabular}
	\end{center}
	Hãy lập bảng thống kê số lượng huy chương đồng, bạc, vàng mà đội tuyển Việt Nam đạt được trong kỳ thi này.
	\loigiai{
		\\
		\begin{tabular}{|c|c|}
			\hline
			\textbf{Loại huy chương} & \textbf{Số lượng} \\ \hline
			Đồng & $2$ \\ \hline
			Bạc & $2$ \\ \hline
			Vàng & $2$ \\ \hline
		\end{tabular}
	}
\end{bt}


%%%%%%%%%%%%%%% BÀI 7
\begin{bt} 
	Tìm $x, y$ biết: $\overline{1xy}$ là bội của $9$ và $x$ là số nguyên tố lẻ nhỏ nhất.
	\loigiai{ $135$}
\end{bt}